% !TeX root = ../papers/nc-p3.tex
\section{楕円曲線の切断環からの反転窓と全次数認証}
\label{sec:elliptic-reversible-pencil}
滑らかな種数1の射影曲線$C$、順序付きの次数4の直線束$L,M$、$N=LM$と
基点を持たない鉛筆$W\in\Gr(2,H^0(N))$を取る。直線束の積はテンソル積を表す。
本節の曲線は元代数の点加群を並べた曲線ではない。反対切断の左右の核にある二次式から復元する。

\begin{theorem}[楕円鉛筆の全ファイバーの正則性]\label{thm:elliptic-pencil-all-as}
一次空間を$E=H^0(L)^*,F=H^0(M)^*$と置く。第一窓の六次元の関係空間を
\[
 R=\ker\bigl(H^0(L)\otimes H^0(M)\longrightarrow H^0(N)/W\bigr)^\perp
 \subset E\otimes F
\]
とし、$R,R^{\rm op}$を交互に繰り返すと標準14641型の左右Koszul AS正則対象を得る。
$L\not\simeq M$なら第三階数は常に8であり、その両側の二次核はそれぞれ
$C\subset\mathbb P(E), C\subset\mathbb P(F)$を定義する楕円四次曲線の二次式の鉛筆である。
$L\simeq M$の場合は四二次式の対称完全交叉のチャートに入る。
\end{theorem}
\begin{proof}
$L_i$を交互に$L,M$とし、$D_{ij}=L_i\cdots L_{j-1}$、$D_{ii}=\mathcal O_C$とする。
次数$i,j$に$H^0(D_{ij})$を置き、切断の積で結んだ代数$S$がKoszulであることを示す。
各始点$i$で$K_0=\mathcal O_C$とし、
\[
 G_n=K_n L_{i+n},\qquad
 0\longrightarrow K_{n+1}\longrightarrow H^0(G_n)\otimes\mathcal O_C
 \longrightarrow G_n\longrightarrow0
\]
を順次作る。次数2以上の直線束は大域生成されて$H^1=0$である。
次数2の束$D$の評価核は$D^{-1}$である。次数3または4では滑らかな点$p$を取り、
$D(-p)$と$D$の評価列を比較する。切断列は全射であり、商の評価は
$\mathcal O_C\to k(p)$なので、$D$の評価核は$D(-p)$の評価核と$\mathcal O(-p)$の拡大である。
よって次数2,3,4の直線束の評価核は次数$-2,-1$の直線束を商に持つフィルトレーションを許す。
この三次数の束でフィルターされた束にも同じことが成立する。
実際、拡大の切断列は$H^1=0$で完全、商と部分束が大域生成されるので中項も大域生成され、
三つの評価核も完全列をなす。従って$G_0=L_i$から始め、$K_{n+1}$に次の次数4の束を
掛けると再び次数2,3の直線束でフィルターされる。

上の評価列を$D_{i+n+1,j}$でひねり、大域切断を取る。尾部の長さが1以上なら
$K_{n+1}D_{i+n+1,j}$のフィルトレーションの各商は正次数で$H^1=0$、
長さ0なら評価は大域切断上恒等である。従って始点単純の線型分解が得られる。
具体的に第一射は$H^0(L_i)\otimes e_{i+1}S\to e_iS$、その核の位置$j>i$での
成分は$H^0(K_1D_{i+1,j})$である。以後も同じ議論で、この核を次の位置の
$H^0(G_1)$が生成し、各段の核と生成元が同様に続く。逆順でも成立する。
特に$S$は一次生成二次関係である。

偶奇二頂点に折り返すと、$W$の基底$a,b$はそれぞれ両頂点からの長さ2の切断の和として中心元になる。
$a$は正則で、$b$も$a$を割った後で正則である。実際、$bs=at$なら基点自由性から
$s$は$a$の零因子を重複度込みで含み、$s/a$は所定の尾部の直線束の大域切断となる。
次数が2未満ならその束の次数は負である。
$H=S/(a,b)$はKoszulである。一般にKoszul代数から次数2の中心正則元$z$を割るとき、
環変更Torスペクトル系列の二行は商のTorとその内部次数2の移動である。
商で最初に非対角成分が出るホモロジー次数$p$を取れば、第零行からの唯一あり得る微分の先は
次数$p-2$の第一行にあり、最小性から該当成分は零である。非対角成分が元の代数のTorに残って矛盾する。
この証明は次数零の代数が二頂点の半単純代数でも同じである。
長さ$n>0$の$S$の次元は$4n$なので、二つの次数2の正則元を割った各始点の次元列は$(1,4,6,4,1)$となる。

随伴完全対合も直接得られる。任意の$D=D_{ij}$、$0\le j-i\le4$について基点自由鉛筆の完全列
\[
 0\to DN^{-2}\to (DN^{-1})^{\oplus2}\to D\to0
\]
を使う。$Q(D)=H^0(D)/W H^0(DN^{-1})$と書けば、連結写像とSerre双対により
$Q(D)\simeq Q(N^2 D^{-1})^*$を得る。これは切断の積の後に
$H^0(N^2)/W H^0(N)\simeq H^1(\mathcal O_C)$へのtraceを取る対合である。
従って$H$は有限Frobeniusで、相補次数の対合により二次双対の左右のKoszul単純分解は
四移動して双対になる。大域次元4、左右AS条件と標準Hilbert級数が得られる。
その関係は表示の$R$と反転である。

$\ell$を上記traceとする。標準テンソルは四切断の積に$\ell$を施したものなので、
反対切断は$H^0(L^2)\times H^0(M^2)\to k$の積対合を通る。
次数4の完全線型系による楕円四次曲線は二二次式の完全交叉で、対称二乗は全二次切断に全射である。
この事実は二次式が少なくとも二本あり、非退化既約四次曲線を含む独立二式は共通平面を
持てず、その完全交叉が次数4を使い切ることからも従う。
$L\not\simeq M$なら$L^2N^{-1}$と$M^2N^{-1}$は非自明な次数零の束である。
先の鉛筆完全列を$D=L^2$にも適用すると、八次元の両切断空間がそのまま完全に対合する。
従って階数と核は主張通りである。$L=M$なら切断環は通常の楕円四次曲線の環で、
二本の定義式に$a,b$の二次式持ち上げを加えた対称完全交叉になる。
\end{proof}

\begin{theorem}[曲線データの復元と被覆]\label{thm:elliptic-pencil-reconstruction}
$L\not\simeq M$の上記全対象では、添字保持同型の必要十分条件は曲線同型が順序付きの
$(L,M,W)$を対応させることである。特に最小周期1の条件は$(M,L,W)$との曲線同型である。
このチャートは定理\ref{thm:reversible-component}の44次元成分の稠密開部分を覆う。
\end{theorem}
\begin{proof}
まず添字保持同型に現れる反対二位置の基底変換を揃えられることを示す。
第三切断の像は対称二乗内の八次元部分である。四次元空間の非スカラーな$g$に対し、
対称形式$A$で$A(gv,w)=A(v,gw)$を満たす空間の次元は7以下である。
実際、$v,gv$が独立な$v$を固定すると、$w\mapsto (gv)w-v(gw)$は対称二乗への階数3の写像である。
$v$を法として見ると核は$kv$に限られる。従って対称二乗の八次元部分を対称二乗に移す
二基底変換はスカラーを除いて同一である。これを両側に適用すれば、標準テンソルの同型は
$\PGL(E)\times\PGL(F)$で比較でき、両側の二次核から二つの埋込み楕円曲線$C_E,C_F$を得る。

両曲線の間の同定もテンソルから一意に決まる。$C_E\times C_F$上の束
$P=\mathcal O_{C_E}(2)\boxtimes\mathcal O_{C_F}(2)$を取る。
テンソルはその大域切断上の汎関数であり、実際の同定のグラフへ制限した切断に$\ell$を施す。
相異なる二つの同定のグラフ$G,G'$があったとする。$P(-G-G')$は豊富である。
実際、$C_F\to\mathbb P^1$の次数2の射で同定$\phi$の前後の点が等しい像を持つ条件を取ると、
$G_\phi$と対合を合成したグラフとの和が双次数$(2,2)$の外積束の因子となる。
各グラフは差の射のファイバーとしてnefである。従って$P(-G-G')$は次数$(4,4)$の
豊富な外積束に二つのnefなグラフ因子を加えたものとなる。
曲面の標準束は自明なのでKodaira消滅より両グラフの合併への制限は全射である。
二つのグラフの線型包の交わりは、その交点スキームの線型包に含まれる。
実際、合併上の切断は交点で一致する二切断の対であり、両制限は全射である
（各グラフ上の束は次数16で、下記の交点を引いても正次数）。

異なる二同定の一致点は高々四つで、全て横断的である。差を自己同型と見て固定点があればそれを原点に取る。
原点を保つ非自明な楕円自己同型は有限位数$n\ge2$で微分は1ではない。
巡回群による商へのHurwitz公式から固定点数は$2n/(n-1)$以下である。
従って両線型包の交点のテンソルは高々四つの点評価の和であり、反対階数8に達しない。
これでグラフの一意性を得る。グラフ上の混合乗法を$\mu$とすれば
標準テンソルの隣接切断の核を$\mu$で送ったものが$W$である。これは$H$の二次成分の表示と
その完全積対合による。逆に曲線同型は全切断の積と関係を移すので十分である。

楕円曲線の原点を平行移動して$L=\mathcal O(4o)$と置け、残る原点の選択は有限である。
$M$の選択は一次元、鉛筆の選択は$\Gr(2,8)$の非空開集合で12次元である。
曲線の一次元を加えて14次元となる。これらのデータは既約な族をなし、$L=M$にも特殊化できる。
例えば原点付きWeierstrass族上で$M=\mathcal O(3o+p)$と置き、相対Grassmann束を取ればよい。
標数零で$L$を保つ曲線自己同型群は有限であり、上記復元から射影基底を加えた第一窓の像は44次元である。
ここで反転を繰り返す対象では、第一窓の基底変換はそのまま全関係列の基底変換に延び、
第一窓を固定する射影変換も上の復元により有限群をなす。

例の四つの基点自由な二次式から一般の二式を選ぶと、Bertiniの定理で滑らかな楕円四次曲線が得られる
（完全交叉は連結で随伴により種数1）。残る二式はそこで基点自由な鉛筆となる。
従って$L=M$で二十点の第一窓はこの族の像に含まれる。全てのGramテンソルは両因子対称で
隣接階数6なので、その像は前節の二重交代化の小行列式スキーム$Z$に入る。
例を通る局所成分は唯一で次元44であるから、上の既約像の閉包はその成分である。
構成可能な稠密像は成分の稠密開部分を含み、そこにある全ての幾何点が前定理による全次数認証を持つ。
\end{proof}

\paragraph{有理数上の照合。}
\path{computations/verify_elliptic_reversible_pencil.py}では、八点の評価行列$X$のうち後四点を
一平面に取り、$Y$の列を$X^{\mathsf t}$の核の基底に取る。八点を通る二二次式の交叉は滑らかで、
その超平面束を$L$、八点の和を$P_8$、$M=\mathcal O(P_8)L^{-1}$とすると、$Y$は$M$による八点の評価を表す。
実際、$L$の評価の余核の双対は、因子$P_8$の完全列とSerre双対により$H^0(M)$であり、
各行の自明化で留数の非零定数を吸収できる。$N=\mathcal O(P_8)$の切断の評価像は
成分和が零の超平面である。$P_8$を零とする切断と、評価が
$(1,2,3,4,5,6,7,-28)$となる切断で基点自由鉛筆を作ると、traceはこの八係数の逆数を重みに持つ。
これは第二切断と任意の$N$の切断との積を消すことから分かる。両側の四次判別式から得る不変量の一致、
反転窓の零点、二つの切断階数6と8を有理数上で検算する。

\begin{corollary}[周期の判定と一般の直線束に対する一周期層]\label{cor:elliptic-pencil-period}
原点を固定し、次数4の束$L,M$の点和を$l,m\in C$とする。
$L\ne M$の場合、一周期となる必要十分条件は
\[
 u\in\operatorname{Aut}(C,o),\quad (1+u)(l-m)=0,\quad 4a=l-u m
\]
を満たす$\phi(z)=uz+a$の少なくとも一つが$W$を保つことである。
$j\ne0,1728$かつ$2(l-m)\ne0$なら、固定した$(C,L,M)$の鉛筆の12次元空間で
一周期部分の次元は6である。
\end{corollary}
\begin{proof}
$\phi^*$は次数4の点和$z$に$u^{-1}(z-4a)$として作用するので前定理で判定式を得る。
表示した一般条件では$u=-1$のみで、$a$の選択は十六個である。
この対合による商を$\pi:C\to\mathbb P^1$とすると、
$N=L\otimes\phi^*L$はノルムの引き戻し$\pi^*\mathcal O(4)$である。
$\pi_*\mathcal O_C=\mathcal O\oplus\mathcal O(-2)$より切断の固有空間次元は5と3である。
二次元不変部分空間の正負次元が$(2,0)$なら$\Gr(2,5)$の基点自由な非空開部分で次元6である。
$(1,1)$なら$\mathbb P^4\times\mathbb P^2$の開部分で、一般の正切断は負切断の有限個の零点を避ける。
$(0,2)$なら分岐上の共通零点を持つ。有限個の対合の和を取っても次元は6である。
\end{proof}

\begin{proposition}[正次元の点対応を持つ楕円鉛筆]\label{prop:elliptic-pencil-dihedral}
位数4の点$t\in C$と$T=\langle t\rangle$を取り、$f:C\to(C/T)/\{\pm1\}=\mathbb P^1$を取る。
$L=\mathcal O(3o+[t])$, $M=\mathcal O(3o+[-t])$とし、
$W=f^*H^0(\mathcal O(1))$とする。この標準AS対象は一周期、第三階数8である。
第一窓の点対応の被約支持は、正規化が$C$である一つの曲線と四つの孤立被約点からなる。
\end{proposition}
\begin{proof}
一般の$f$のファイバーは$p+T$と$-p+T$の和で、八点の点和は零であるから
$f^*\mathcal O(1)\simeq LM$である。$L,M$は異なり、反転がこれらを交換して鉛筆を保つ。
零点対とは$f$のファイバー因子の順序付き分割$D+D'$で$D\in|L|,D'\in|M|$となるものにほかならない。
これは第一関係の直交空間が$\mu^{-1}W$であることによる。

$p+T$から$k$点を取ると$D$の点和は$(2k-4)p$に$T$の元を加えたものである。
一般点で可能なのは$k=2$のみである。二点ずつの添字を$I,J\subset\mathbb Z/4$とすると条件は
$\sum I+\sum J=1\pmod4$である。得られる八組は
$(I,J)=(\{0,1\},\{1,3\})$から$(I,J)\mapsto(I+1,J-1)$と交換により
安定化群なしで一軌道となる。従ってこの曲線は$C$から次数1で被覆され、他の支配曲線はない。
各ファイバーの分割は有限個なので垂直な曲線もない。

分岐外で$k=1,3$なら$2p\in T$となって分岐外に反する。$k=4$の分割は
$D=p+T$、$4p=-t$となり、$T$の平行移動を除いて四つある。$k=0$は$p$を反転した同じ分割を与える。
これらは分岐外にある。単純な八点を局所的に順序付けると、分割の枝は局所的なファイバーパラメータで
表され、$|L|$に入る式$4p=-t$の微分は非零である。従って四点はいずれも孤立被約である。
最後に$2p=ht\in T$ならファイバーは$p+T$の各点を二回ずつ含む。
$D$の多重度が$(2,2,0,0)$または$(1,1,1,1)$の置換なら添字和が偶数で、
$4p=2ht$を加えても点和は$t$にならない。残る$(2,1,1,0)$型では重複点を両組に一回ずつ、
二つの単純点を別々の組に入れて変形できる。第二組の添字を$h$だけ戻して$-p+T$で書けば、
点和の条件がそのまま$\sum I+\sum J=1$になるので上の曲線の極限に属する。
分岐点での埋め込み成分の有無はこの被約支持の記述には用いない。
\end{proof}

\begin{theorem}[節点・尖点の整な境界]\label{thm:elliptic-pencil-integral-boundary}
$C$を節点または尖点を持つ整な算術種数1のGorenstein曲線とし、$L,M$は可逆層とする。
同じ次数4の二束と基点自由鉛筆の構成では、全ファイバーの標準AS正則性、
$L\not\simeq M$での第三階数8、および曲線同型による添字保持同型の必要十分条件が保たれる。
$L\not\simeq M$の節点層と尖点層の同型類の族はそれぞれ13次元、12次元であり、
上の滑らかな楕円チャートの特殊化に入る。
\end{theorem}
\begin{proof}
双対化層は自明である（次数零で非零切断を持つ）。長さ$h$の部分スキーム$Z$について
$\mathcal I_Z^*$は次数$h$の階数1無捩れ層である。実際、
$0\to\mathcal I_Z\to\mathcal O\to\mathcal O_Z\to0$を双対化すると、
局所Gorenstein双対性により最後の$\mathcal E xt^1(\mathcal O_Z,\mathcal O)$の長さが$h$である。
$\deg D>h$なら$D\to\mathcal I_Z^*$の非零写像は次数により存在しない。
従ってSerre双対で$H^1(D\mathcal I_Z)=0$となり、
次数2以上の束は大域生成、次数3以上では長さ2の部分スキームを分離する。
前定理の評価核のフィルトレーションでは、束の次数を下げる点を滑らかな点に選べば全証明が通る。
鉛筆の列、二中心元の正則性、完全対合の証明もそのまま通る。
非自明な次数零の束は切断を持たず、次数4の埋込みは同じ次数の議論で二二次式完全交叉なので、
反対階数と両側の曲線を再び復元できる。

同定の一意性を正規化$\nu:\mathbb P^1\to C$上で示す。各側の八次元二次切断空間から、
導手イデアルを掛けた七次元部分空間を取る。正規化では、節点の両原像または尖点の二重原像で
消える切断全体であり、固定因子を除けば$H^0(\mathbb P^1,\mathcal O(6))$である。
これは節点の二枝での値の貼り合わせ、尖点での一次係数の貼り合わせからも従う。
八次元完全対合を両側で余次元1の部分へ制限すると階数は6以上である。
この制限テンソルは、各同定が与える正規化のグラフ上の$(6,6)$切断を通って定まる。
実際、正規化上のグラフで消える切断は固定因子を戻した後でも元の整な曲線上で零である。
異なる射影直線の同定二つのグラフは長さ2で交わる。
合併のイデアルでひねると$\mathcal O(4,4)$となり$H^1=0$だから、
両グラフの共通線型包にあるテンソルの行列階数は高々2である。
非被約な長さ2の交点でも、積の対合はその長さ2の座標環を通る。従って同定は一つしかない。
正規化で一致する射は整な曲線上でも一致し、$W$は前と同じ隣接核から戻る。

抽象節点曲線は一種で、滑らかな部分は乗法群である。原点を保つ群は反転を含む有限群で、
平行移動で$L$を規格化した後は$M$の一次元と鉛筆の12次元が残る。
尖点の場合、滑らかな部分は加法群で、さらに原点を保つ乗法群が作用する。
次数4の点和を$L$では0、$M$では1に規格化できるので鉛筆の12次元だけが残る。
これらの自己同型の記述は、正規化上で節点の二原像を保つ群、尖点の一原像を保つアフィン群から得られる。
特殊化については$|L|$の二二次式を一般の二式の方へ変形すれば滑らかな完全交叉へ移る。
相対超平面束で$L$を延長し、$M$は特異点を避ける一般の四点の因子で表示して、
その滑らかな四点をエタール局所的に持ち上げる。$H^1(N)=0$なので$W$も延長でき、基点自由性は開である。
隣接関係は階数一定の乗法の核の直交空間であり、これが求める特殊化を与える。
\end{proof}

\begin{proposition}[有理曲線の環状境界]\label{prop:elliptic-pencil-polygons}
$C$が二、三または四本の滑らかな有理曲線を節点で環状につないだ曲線
（二本なら異なる二点でつなぐ）とする。$L,M$が各成分で同じ正の多重次数$d$を持ち、
$\sum d_i=4$とする。基点自由な全鉛筆について先と同じ構成は標準AS正則で、
$L\not\simeq M$なら第三階数8となり、両側の二次核はこの環状四次曲線を切り出す。
これらは滑らかな族の特殊化に入る。$L\not\simeq M$の場合には曲線データによる同型判定も保たれ、固定多重次数での有理商の次元は、成分数$n=2,3,4$に対して$14-n$である。
\end{proposition}
\begin{proof}
環状曲線では、正規化の各成分の切断を二つの端点で貼り合わせる。双対化層は自明である
（各成分に、両端だけで極を持ち留数が逆となる対数微分を置く）。
非負の多重次数で総次数が正の束は$H^1=0$となる。双対の切断が負次数成分で消えれば
次数零成分の列にもその消滅が伝わるためである。総次数が2以上なら大域生成される。
実際、正次数成分が二つ以上あれば各成分上で二つの端点値を独立に指定でき、
次数零の鎖では値を定数倍で引き継ぐ。正次数成分が一つだけなら、
その両端値に一つの比例条件を課した次数2以上の多項式系になり、基点を持たない。

評価核の証明では多重次数が$0\le b\le d$の束を取る。総次数3以上の場合は
正次数成分から滑らかな点を引いて非負の総次数2まで下げる。
最後の鉛筆の逆束に次の$L_i$を掛けると、非負の多重次数$d-b'$、総次数2となる。
他の評価核の商$\mathcal O(-p)$に掛けたものも総次数3で同じ範囲に残る。
これで線型分解の帰納と、尾部をひねった核の$H^1$消滅が成立する。
鉛筆の一般の第一切断$a$はどの成分でも恒等的に零にならず、第二切断は$a$の零点で単元である。
従って正則列と双対性の証明も通る。非自明な多重次数零の束には大域切断がない。

正の多重次数$d$の完全線型系は各成分を有理正規曲線に写す。
座標は各節点の値と、各成分上で両端で消える$d_i-1$個の内部座標である。
各成分は両端と自分の内部座標が張る空間に入り、異なる成分は指定の節点でだけ横断的に交わる。
二成分の場合も、二つの二次曲線を含む両平面の共通直線、または三次の有理正規曲線の割線は、
非線型成分と両端でのみ交わる。この表示で次数3の部分曲線が一平面に入ることはない。
独立な二本の消滅二次式は次元から存在し、もし共通平面因子があればその外の成分が
残る二平面の共通直線に入るので次数3以上が共通平面に残ってしまう。
従って二式は完全交叉を定め、次数と算術種数からこの曲線に等しい。
この完全交叉と$M$の一般の滑らかな四点を動かせば、前定理と同様に平滑化できる。

同定の復元を示す。一般の$a\in W$の零点は成分$i$に$r_i=2d_i$個ずつあり、
全て単純で節点を避ける。traceはこの点評価の非零重み付き和である。
実際、鉛筆列の連結写像はもう一方の切断$b$を選んで$a$の零点で$1/(ab)$の留数を取る写像であり、
$a$の単純性、$b$の非消滅、双対化層の生成微分から各重みは非零である。
$2L$を意味する束$L^2$の切断双対で、成分$i$の評価の線型包を$P_i$とする。
$M^2$側の同様の空間を$Q_i$とし、実際のグラフで対応するよう番号を付ける。
どちらも次元$r_i+1$である。節点評価と成分内部の$r_i-1$座標による基底では、
各$P_i,Q_i$は対応する座標部分空間となる。traceテンソルは$T=\sum T_i$で、
$T_i\in P_i\otimes Q_i$は上の$r_i$個の評価から作る階数$r_i$のテンソルである。

片方で成分$i$だけを台とする切断は端点二つで消え、固定因子を除けば
$H^0(\mathcal O(r_i-2))$である。その縮約で得る像は$Q_i$内で次元$r_i-1$を持つ。
別の対応成分を仮定すれば、この像は二成分の線型包の共通部分に入る。
$d_i\ge2$では共通部分は高々二次元なので矛盾する。
$d_i=1$で三成分以上なら、異なる成分の共通部分は高々一節点の直線である。
しかし縮約の像は二つの非節点での非零重み付き評価の和であり、
二次Veronese曲線の異なる三点の独立性から節点の直線には入らない。
二成分で$d_i=1$なら他方は$d_j=3$であり、次数でも区別される。従って成分対応は一意である。

三成分以上の場合は成分対応で節点対応も定まる。二つのグラフ表示で$T_i$が異なるとしても、
その差は他の$P_j\otimes Q_j$の和に入る。座標基底で共通するのは両端に対応する
二本の節点評価テンソルだけであり、これらは双方のグラフの線型包に入る。
従って$T_i$自身が二つの成分グラフの線型包に入る。
異なる同定なら二グラフの交点はその両端二点であり、共通線型包は二つの節点評価の包に等しい
（合併のイデアルでひねると$\mathcal O(r_i-2,r_i-2)$となる）。
$r_i>2$なら階数に反し、$r_i=2$でも、二つの非節点の評価が張る割線は節点二つの割線と異なる。

二成分の場合には少なくとも一方で$d_i\ge2$となる。片側だけをその成分支持に制限した
階数$r_i-1\ge3$のテンソルは正規化上の$(r_i-2,r_i)$グラフを通る。
二同定に共通ならランクは2以下となる（合併でひねると$\mathcal O(r_i-4,r_i-2)$で$H^1=0$）。
従ってこちらの同定と両節点の対応が定まる。
同じグラフ上の二つの$T_i$の差は、両側で両節点の評価が張る二次元空間のテンソル積に入る。
この四次元空間のうちグラフの線型包に入るのは、同じ節点どうしの二項だけである。
実際、両節点を$0,\infty$と取れば、異なる両端の単項式対の積は、
両端で消える内部単項式の対の積としても書ける。後者での値は零なので交差する二係数は零である。
ゆえに残る成分の$T_j$は二つの節点項を除いて定まり、先の割線の議論でその同定も一意となる。
$W$を戻せば同型の必要十分性が従う。

最後に、各束の同じ多重次数での貼り合わせパラメータは一つの乗法群である。
各有理成分の両端を保つ群の連結部分は$n$次元トーラスで、一般の鉛筆の安定化群は有限となる。
三成分以上なら一般の鉛筆は異なる節点を異なる値に送るため、基底の射影変換は有限個である。
これは節点値を独立に指定できることから従う。二成分なら少なくとも一成分の射は次数4以上で、
一般には三つ以上の分岐値を持つ（その成分への切断の制限は全射）。
基底の射影変換を固定した後で各曲線上の有限射の持ち上げも有限個である。
従って有理商の次元は$1+1+12-n$となる。
\end{proof}

\paragraph{異なる多重次数の境界での第三切断。}
環状曲線上の$L,M$が相異なる正の多重次数を持つ場合も、鉛筆のtraceテンソルの
反対切断階数だけは全て決まる。両端の対称二乗は各々の全二次切断へ全射なので、
鉛筆完全列による完全対合からその階数は
\[
 8-2h^0(LM^{-1})
\]
である。実際、$W\otimes H^0(LM^{-1})$からの積の核は
$H^0(M^{-2})=0$である。二成分で多重次数が$(3,1),(1,3)$（逆順も可）なら、
差の束の切断は一成分上の次数2の多項式で両端で消えるものからなり、階数は6である。
他の異なる正の多重次数では差は$+1,-1$と零だけである。負次数成分で切断が消え、
両側から零次数の鎖を経て正次数成分の両端で消えるため$h^0=0$、階数は8である。

\begin{corollary}[正の多重次数の必要十分条件]\label{cor:elliptic-polygon-degree-balance}
同じ環状曲線上の任意の正の多重次数$d,e$、$\sum d_i=\sum e_i=4$を許し、
第一窓を定理\ref{thm:elliptic-pencil-all-as}の表示で定める。鉛筆が基点自由なら、
反転二窓の代数が標準Koszul AS正則となる必要十分条件は$d=e$である。
異なる場合は全鉛筆で指数成長する。
\end{corollary}
\begin{proof}
十分性は命題\ref{prop:elliptic-pencil-polygons}である。
相補的な空でない二つの連続部分鎖$I,J$を取る。$I$だけを台に持つ$L$の切断は境界二点で消え、
次元$d(I)-1$を持つ。実際、正次数の各成分では両端値を独立に指定でき、内部節点で貼り合わせると
鎖上の切断次元は$d(I)+1$となり、両端値の零条件が二つである。
同様に$J$に支えを持つ$M$の切断の次元は$e(J)-1$である。
この二空間の積は順逆とも零である。従って元代数の矢をそれぞれの双対へ射影すると、
積の核の直交空間である全二次関係の像が消える。
この経路代数への全射により、長さ$2m$の次元は少なくとも
$((d(I)-1)(e(J)-1))^m$となる。

二成分で$d\ne e$なら、成分の順序と両束の役割を選んで一方の束は一成分に次数3、
もう一方は残る成分に次数2以上を持つ。三成分で異なる場合は、$J$を$M$の次数2の成分、
$I$を残り二成分とすれば$d(I)=3$である。四成分では正の多重次数は一種類しかない。
従って異なる全場合で指数成長し、標準Hilbert級数に反する。
\end{proof}
